\subsection{SIMPLICITY THEOREMS}

Lemma 2. Let H be a normal subgroup of G. There exists a subset X of S such that BH = G\_X and such that every element of X commutes with every element of S-X.

Since BH is a subgroup of G containing B, there exists a unique subset X of S such that \(BH = G_{X}\) (Th. 3).

Let \(s_1 \in X\) and \(s_2 \in S - X\) ; let \(n_1\) and \(n_2\) be representatives in \(N\) of \(s_1\) and \(s_2\) , respectively. Then \(n_1 \in G_X = BH\) and there exists \(b \in B\) such that \(bn_1 \in H\) . Since \(H\) is normal in \(G\) , the element \(h = n_2bn_1n_2^{-1}\) of \(G\) belongs to \(H\) . This means that

\[
h \in \mathrm{C} (s _ {2}). \mathrm{C} (s _ {1}). \mathrm{C} (s _ {2}).
\]

If the length of \(s_2s_1s_2\) is equal to 3, Cor. 1 of Th. 2 implies that

\[
\mathrm{C} (s _ {2}). \mathrm{C} (s _ {1}). \mathrm{C} (s _ {2}) = \mathrm{C} (s _ {2} s _ {1} s _ {2}),
\]

and hence that \(h \in H \cap C(s_{2}s_{1}s_{2})\) . Since \(H \cap C(s_{2}s_{1}s_{2})\) is non-empty, \(s_{2}s_{1}s_{2} \in W_{X}\) . As \((s_{2}, s_{1}, s_{2})\) is a reduced decomposition, it follows that \(s_{2} \in X\) , contrary to our assumption.

Thus \(l_{\mathrm{S}}(s_2s_1s_2) \leqslant 2\) ; if \(l_{\mathrm{S}}(s_2s_1s_2) = 1\) , then \(s_1s_2 \in \mathrm{S}\) and so \((s_1s_2)^2 = 1\) , or \(s_1s_2 = s_2s_1\) . If \(l_{\mathrm{S}}(s_2s_1s_2) = 2\) , property (E) of §1, no. 5 implies that \(s_2s_1 = s_1s_2\) , since \(s_1 \neq s_2\) . Q.E.D.

The following property of a group U enters into Th. 5 below:

\begin{enumerate}
\def\labelenumi{(\Alph{enumi})}
\setcounter{enumi}{17}
\tightlist
\item
  For any normal subgroup V of U distinct from U, the commutator subgroup (cf.~Algebra, Chap. I, §6, no. 8) of U/V is distinct from U/V.
\end{enumerate}

Every soluble group satisfies (R); in particular, every abelian group satisfies (R). It can be shown that the symmetric group \(S_{n}\) satisfies (R) (cf.~Exerc. 29).

THEOREM 5. Let Z be the intersection of the conjugates of B, let U be a subgroup of B and let \(G_{1}\) be the subgroup generated by the conjugates of U in G. We make the following assumptions:

\begin{enumerate}
\def\labelenumi{(\arabic{enumi})}
\item
  U is normal in B and B = UT.
\item
  U has property (R.).
\item
  \(\mathbf{G}_1\) is equal to its commutator subgroup.
\item
  The Coxeter system (W, S) is irreducible (cf.~§ 1, no. 9).
\end{enumerate}

Then every subgroup H of G normalised by \(G_{1}\) is either contained in Z or contains \(G_{1}\) .

First we show that \(G = G_{1}T\) . The group \(G_{1}T\) contains B and hence is its own normaliser (Th. 4); but as N normalises \(G_{1}\) and T, it also normalises \(G_{1}T\) , so \(N \subset G_{1}T\) . Since G is generated by B and N, it follows that \(G = G_{1}T\) .

Next, set

\[
\begin{array}{c} \mathrm{G} ^ {\prime} = \mathrm{G} _ {1} \mathrm{H}, \quad \mathrm{B} ^ {\prime} = \mathrm{B} \cap \mathrm{G} ^ {\prime}, \quad \mathrm{N} ^ {\prime} = \mathrm{N} \cap \mathrm{G} ^ {\prime}, \\ \mathrm{T} ^ {\prime} = \mathrm{T} \cap \mathrm{G} ^ {\prime} = \mathrm{B} ^ {\prime} \cap \mathrm{N} ^ {\prime} \quad \text {and} \quad \mathrm{W} ^ {\prime} = \mathrm{N} ^ {\prime} / \mathrm{T} ^ {\prime}. \end{array}
\]

We have \(G = G'T\) since \(G'\) contains \(G_{1}\) , and hence \(N = N'T\) . The inclusion of \(N'\) into N thus defines, on passing to the quotient, an isomorphism \(\alpha : W' \to W\) . Let \(S' = \alpha^{-1}(S)\) .

We now show that \((\mathrm{G}^{\prime},\mathrm{B}^{\prime},\mathrm{N}^{\prime},\mathrm{S}^{\prime})\) is a Tits system. Since G = BNB and B = TU = UT, we have G = UNU. Since U is a subgroup of \(G^{\prime}\) , it follows that \(G^{\prime} = UN^{\prime}U\) ; since \(U \subset B^{\prime}\) , this proves (T1). Axiom (T2) is satisfied since \(\alpha\) is an isomorphism. Let \(w \in W\) and let \(w' = \alpha^{-1}(w)\) be the corresponding element of \(W^{\prime}\) . We have

\[
\mathrm{B} w \mathrm{B} = \mathrm{B} w \mathrm{B} ^ {\prime} = \mathrm{B} w ^ {\prime} \mathrm{B} ^ {\prime}, \quad \text { since } \mathrm{B} = \mathrm{B} ^ {\prime} \mathrm{T}.
\]

From this we conclude that \(G' \cap BwB = B'w'B'\) , which means that the inclusion of \(G'\) into G defines on passing to the quotient a bijection from

\(B'\backslash G'/B'\) to \(B\backslash G/B\) . Axiom (T3) follows immediately. Axiom (T4) follows from \(B = B'T\) .

The subgroup H is normal in \(G'\) . By Lemma 2 applied to \((G', B', N', S')\) , there exists a subset \(X'\) of \(S'\) such that \(B'H = G_{X'}'\) and every element of \(S' - X'\) commutes with every element of \(X'\) . In view of assumption (4), there are only two possibilities:

\begin{enumerate}
\def\labelenumi{\alph{enumi})}
\item
  \(X' = \varnothing\) , i.e.~\(B'H = B'\) , so \(H \subset B' \subset B\) . If \(g \in G\) , then \(g = g_1 t\) with \(g_1 \in G_1\) , \(t \in T\) , and \(H \subset g_1 Bg_1^{-1}\) since \(G_1\) normalises \(H\) . Thus \(H \subset gBg^{-1}\) , and since \(Z\) is the intersection of the \(gBg^{-1}\) , we have \(H \subset Z\) .
\item
  \(X' = S'\) , i.e.~\(B'H = G'\) . Since \(G = G'T\) , we have
\end{enumerate}

\[
\mathrm{G} = \mathrm{B} ^ {\prime} \mathrm{HT} = \mathrm{HB} ^ {\prime} \mathrm{T} = \mathrm{HB}.
\]

As B normalises U, every conjugate of U is of the form \(hUh^{-1}\) with \(h \in H\) . Such a subgroup is contained in the group UH, hence \(G_{1} \subset UH\) by the definition of \(G_{1}\) . Thus, we have the isomorphisms

\[
\mathrm{U} / (\mathrm{U} \cap \mathrm{H}) \cong \mathrm{UH} / \mathrm{H} = \mathrm{G} _ {1} \mathrm{H} / \mathrm{H} \cong \mathrm{G} _ {1} / (\mathrm{G} _ {1} \cap \mathrm{H}).
\]

By assumption (3), \(G_{1}/(G_{1} \cap H)\) is equal to its commutator subgroup. Assumption (2) now shows that the group \(U/(U \cap H)\) , which is isomorphic to \(G_{1}/(G_{1} \cap H)\) , reduces to the identity element. Hence \(G_{1} \cap H = G_{1}\) and \(G_{1} \subset H\) , which completes the proof.

COROLLARY. Under the assumptions of Th. 5, the group \(\mathrm{G}_{1}/(\mathrm{G}_{1} \cap \mathrm{Z})\) is either simple non-abelian or reduces to the identity element.

Th. 5 shows that \(G_{1}/(G_{1} \cap Z)\) is simple or reduces to the identity element. On the other hand, assumption (3) implies that it is equal to its commutator subgroup. Hence the corollary.

Remarks. 1) Assumptions (2), (3), (4) were not used in the proof that \((\mathbf{G}', \mathbf{B}', \mathbf{N}', \mathbf{S}')\) is a Tits system.

\begin{enumerate}
\def\labelenumi{\arabic{enumi})}
\setcounter{enumi}{1}
\item
  Suppose that \(\mathbf{Z} \cap \mathbf{U} = \{1\}\) . Since \(\mathbf{Z}\) and \(\mathbf{U}\) are normal in \(\mathbf{B}\) , it follows that every element of \(\mathbf{Z}\) commutes with every element of \(\mathbf{U}\) , and hence with every element of \(\mathbf{G}_1\) . In view of the preceding corollary, it follows that \(\mathbf{G}_1 \cap \mathbf{Z}\) is the centre of \(\mathbf{G}_1\) .
\item
  Assumption (3) is implied by the following condition:
\end{enumerate}

(3') U is generated by the commutators \(b^{-1}u^{-1}bu\) with \(u \in \mathbf{U}\) and \(b \in \mathbf{B} \cap \mathbf{G}_1\) .

Examples. 1) Let \(k\) be a field, \(n\) an integer \(\geqslant 0\) , \(G = \mathbf{GL}(n, k)\) , and let \((G, B, N, S)\) be the Tits system described in no. 2. Let \(U\) be the strictly upper triangular subgroup of \(G\) , i.e.~the subgroup of \(B\) consisting of the matrices whose diagonal entries are equal to 1. Condition (1) in Th. 5 is immediate, and so is (2) since U is soluble. Condition (4) is satisfied if \(n \geqslant 2\) . One can show (cf.~Algebra, Chap. II, §10, Exerc. 13) that (3) is satisfied if \(n \geqslant 3\) and \(\operatorname{Card}(k) \geqslant 4\) . Under these conditions, we conclude that \(G_{1}/(G_{1} \cap Z)\) is simple and that \(G_{1} \cap Z\) is the centre of \(G_{1}\) (cf.~Remark 2).

When \(k\) is commutative, \(G_1 = \mathbf{SL}(n,k)\) (cf.~Algebra, Chap. III, § 8, no. 9).

*2) Let \(\mathfrak{g}\) be a simple Lie algebra over \(\mathbf{C}\) , and let \(\mathbf{G}\) be the group of inner automorphisms of \(\mathfrak{g}\) (cf.~Chap. III, §6, no. 2, Prop. 2). By using Th. 5, one can show that \(\mathbf{G}\) is simple non-abelian.*

